\section{AnthropicとAccenture、embedded evaluationで提携}
Anthropicは、フロンティアAIの独立評価に向けてAccentureと提携し、社内にembedded evaluatorを置くと発表した。Facultyが評価・レッドチーム、アライメント評価、セーフガード試験を担い、双方は今後5年で各10億ドル以上の能力構築投資を見込むとしている。一方Xでは、評価主体としてコンサル大手を選ぶ妥当性への批判も観測された。
\dailyxsource{Anthropic (@AnthropicAI)}{https://x.com/AnthropicAI/status/2101039819870937247}

\section{Qwen3.8-Omni-Flashを公式発表}
Qwenは、音声・映像の同時理解とツール実行を統合したQwen3.8-Omni-Flashを発表した。1Mコンテキストを備え、エージェント用途を前面に出す。性能値や先行モデル比のコスト削減率はQwen側の自己報告であり、本Daily版では独立再測定を行っていない。Qwen-MM-Pluginsのオープンソース化も告知された。
\dailyxsource{Qwen (@Alibaba\_Qwen)}{https://x.com/Alibaba_Qwen/status/2100785962414702599}

\section{ChatGPTプラグインで複数アカウント接続}
OpenAIのPlugins/MCP担当者は、ChatGPTの大半のプラグインで複数アカウントを接続できるようになったと告知した。仕事用と個人用の文脈を同一会話で扱えることを狙い、開発者側の追加変更は原則不要としている。MCPサーバにprofile toolを追加するとアカウント表示を改善できるとの説明もあった。
\dailyxsource{Max Stoiber (@mxstbr)}{https://x.com/mxstbr/status/2100966048132718786}

\section{Mistral、不正アクセス主張を公式否定}
Mistralは、自社システムへの不正アクセスがあったとする主張を認識し調査した結果、裏付ける証拠はなくシステムは侵害されていないと公式Xで表明した。否定対象となった元の主張者や攻撃手法、対象システムの詳細は、観測された公式投稿では示されていない。
\dailyxsource{Mistral AI (@MistralAI)}{https://x.com/MistralAI/status/2100824200126529725}

\section{Claudeを用いたOpenAI関連セキュリティ研究の報道}
複数の二次報道では、Hacktron AIの研究者がClaudeを利用し、OpenAIのDiscourseコミュニティの欠陥から従業員ChatGPTアカウントや内部GitHubへ到達したと伝えられた。6,500ドルのバグバウンティ支払いも繰り返し引用されたが、観測窓ではOpenAIまたはAnthropicの公式声明を確認できておらず、報道段階の情報として扱う。
\dailyxsource{Cointelegraph (@Cointelegraph)}{https://x.com/Cointelegraph/status/2100774325481513463}

\section{CNN、米軍AI情報報告をめぐる事案を報道}
CNNは関係者4人の話として、SOCOM分析官がチャットボットを用いて作成した中国船に関する情報報告が誤っており、米軍が拿捕準備まで進んだ後に誤りが判明したと報じた。出来事自体は今春とされ、観測窓内なのは報道とその議論である。使用モデルは特定されていない。
\dailyxsource{MTS (@MTSlive)}{https://x.com/MTSlive/status/2101000189775609861}

\section{TypeSafe Jev周辺で決定モデルの再現が活発化}
OpenRouterなどで、Jevをyes/no、選択、スコアを返す決定モデルとして捉える解説が広がった。Jared PalmerはQwen2.5-0.5BベースのKev-0.5Bを公開し、ArcherはQwen3.8 27Bベースのオープンウェイト再現を予告した。性能やコストの比較値には自己評価が含まれ、独立検証は未了である。
\dailyxsource{OpenRouter (@OpenRouter)}{https://x.com/OpenRouter/status/2101061688338575739}

\section{Cactus Compute、Needle 3を公開}
Cactus Computeは、8--29MBでスライス可能とする自動化向け基盤モデルNeedle 3を公開した。チャットではなく関数呼び出しや型付きレコードを返し、Raspberry Pi 5上で最大4k tok/sのデコードを主張する。一方、第三者比較ではJevなどより品質・速度が劣るとの報告もあり、実装条件を含めた再現確認が必要である。
\dailyxsource{Cactus Compute (@cactuscompute)}{https://x.com/cactuscompute/status/2100685924401295764}

\section{Muse for Macをロールアウト}
MuseはMac版の提供開始を発表した。ダウンロード整理、ファイル探索、メッセージやメモの要約などを、明示的な許可のもとで実行するパーソナルエージェントとして説明している。MetaのAI公式アカウントも同投稿を引用したが、MuseとMetaの組織・資本関係は今回のGrok収集では確定していない。
\dailyxsource{Muse (@Muse)}{https://x.com/Muse/status/2100714397337264444}

\section{NVIDIA、プラットフォーム論と医用AI事例を紹介}
NVIDIAはAll In SummitでのJensen Huangのプラットフォーム論を動画で紹介した。別投稿では、Children's Hospital of PhiladelphiaがMONAIを用い、既存CT・MRI・3D超音波から小児心臓の精密モデルを短時間で生成できるとする事例を取り上げた。所要時間などの数値はNVIDIA広報による説明である。
\dailyxsource{NVIDIA (@nvidia)}{https://x.com/nvidia/status/2100979179655200951}

\section{Bolt.new、DeepSeek V4.1 FlashをForgeに追加}
Bolt.newはDeepSeek V4.1 FlashをBolt Forgeへ追加したと発表した。DeepSeek V4 Proの10倍の利用があった、見た目品質で好評だったとの数値・評価はBolt側の自己報告である。観測窓ではDeepSeek側からの同時公式発表は確認されていない。
\dailyxsource{bolt.new (@boltdotnew)}{https://x.com/boltdotnew/status/2100988806950252683}

\section{Anthropicの11月IPO観測が拡散}
日本語の投資速報アカウントが、Anthropic幹部による投資家面談や11月IPO計画、将来売上見通しについてWSJやNYTの報道として紹介した。ただし観測窓内にAnthropic公式のIPO発表は確認されておらず、原記事本文もGrok収集では確認できていないため、未検証の報道段階として扱う。
\dailyxsource{NABU (@nabu\_aia)}{https://x.com/nabu_aia/status/2101068678884372892}

\section{長ホライズン向けSKILL.state論文が再注目}
長い会話履歴の追記をSKILL.stateへ置き換えるGoogle論文について、長ホライズン技能実行でコンテキスト汚染を抑え、100ターンで累積トークンを16倍削減したとする解説が再共有された。論文そのものや親投稿は観測窓外の可能性があり、正式タイトルやarXiv IDも今回の収集では未確定である。
\dailyxsource{beamnxw (@beamnxw)}{https://x.com/beamnxw/status/2101018409802580300}

\section{GLiClassをJevの先行OSS類似として再提示}
Knowledgatorは、2024年に開発した分類モデルGLiClassをJevのオープンソース類似物として改めて紹介した。DeBERTa、ModernBERT、Qwen系の変種と、focal loss、PPO、LoRAを組み合わせる学習構成を説明している。リポジトリや論文の安定URLは今回のGrok収集では展開されていない。
\dailyxsource{Knowledgator (@knowledgator)}{https://x.com/knowledgator/status/2101023489352228915}
